\section{Related Work}\label{sec:related}

\paragraph{KG-augmented retrieval.}
GraphRAG~\citep{edge2024graphrag} introduced a community-detection
approach in which a KG is constructed and then summarized at
multiple levels of abstraction; query-time retrieval selects
community summaries rather than raw triples.
HippoRAG~\citep{gutierrez2024hipporag} treats retrieval as
personalized PageRank over a KG built from a corpus, mimicking
hippocampal indexing for multi-hop QA. Both approaches construct
the KG \emph{from} the corpus; we instead start from a
production-curated KG and ask how its \emph{surface form} affects
LLM downstream accuracy. The two strands of work are complementary:
GraphRAG/HippoRAG focus on \emph{retrieval} from KGs, while we
focus on \emph{serialization} of an already-retrieved (or
in-context-feasible) KG to the LLM.

\paragraph{KG verbalization and seq2seq.}
KGT5~\citep{saxena2022kgt5} casts KG completion and KGQA as a
text-to-text task, demonstrating that LLMs trained on linearized
KGs can perform multi-hop reasoning. Our setting differs in two
ways: we use general-purpose chat LLMs (no KG-specific training)
and we evaluate \emph{multiple} surface forms head-to-head on the
same facts.

\paragraph{Hyper-relational KGs.}
StarE~\citep{galkin2020stare} provides a message-passing approach
for hyper-relational KGs, demonstrating that qualifier-augmented
relations carry information that binary-triple embeddings lose.
Earlier foundational work includes m-TransH~\citep{wen2016mtransh}
and NaLP~\citep{guan2019nalp} on n-ary relational data. These
papers establish that hyper-relational structure is real and
exploitable for embedding-based tasks. Our finding does not
contradict this: in-context LLMs handle the same hyper-relational
structure (C-TQL preserves role players verbatim) just as well
as flatter forms; the issue is not whether structure can be
exploited but whether the LLM-facing surface form is the right
place to expose it.

\paragraph{In-context behavior of LLMs.}
\emph{Lost in the middle}~\citep{liu2023lostmiddle} documents
LLMs' positional sensitivity within long contexts. While our
results do not directly engage this finding (our dumps fit
comfortably under 32K characters), it reinforces the case that
in-context surface form interacts with model behavior in
non-obvious ways. Surface-form effects in retrieval-augmented
settings are otherwise under-studied; we are not aware of prior
systematic comparisons of TQL/Cypher/NL/JSON-LD on the same
underlying facts.

\paragraph{Closed-world reasoning.}
The closed-world assumption \citep{reiter1978cwa} is foundational
in knowledge representation. Our finding that in-context KG dumps
strip closed-world semantics is, in spirit, a re-statement of
Reiter's distinction in the LLM era: \emph{absence in the dump
does not entail absence in the world}, and the LLM has no syntactic
cue to make the closure. To our knowledge, this paper is the first
to empirically measure the resulting accuracy gap between formal
KG dumps and natural-language summaries on the same out-of-scope
queries.

\paragraph{KG surveys and TypeDB.}
Hogan et al.~\citep{hogan2021kgsurvey} provide a comprehensive
survey of knowledge graphs that situates hyper-relational
representation within broader KR. TypeDB v3's TypeQL language is
documented industrially~\citep{vaticle2023typeql}; to our
knowledge, no peer-reviewed paper has evaluated TypeDB v3 in a
production LLM-agent deployment.
